\chapter{Material de trabajo y planteamiento experimental}

\section{Funcion de este capitulo}

Este capitulo debe situar al lector antes de entrar en la metodologia formal y antes de describir el flujo del sistema. Su papel es explicar con que material se va a trabajar, que tipo de consultas se quieren resolver y por que ese conjunto resulta adecuado para tensionar el nuevo TFM.

\section{Origen y naturaleza del material de trabajo}

\textit{Completar aqui el origen de los datos, el tipo de instrumentos incluidos, la razon de su seleccion y la forma en que se obtendran o congelaran para el proyecto. Conviene mantener esta seccion breve y muy clara.}

\section{Preguntas utiles para orientar este capitulo}

\begin{enumerate}
    \item ?Que activos o instrumentos quieres incluir en el nuevo TFM y por que esos y no otros?
    \item ?Quieres priorizar variedad de clases de activo, realismo de las consultas o facilidad de reproducir el experimento?
    \item ?Que papel juegan los datos historicos descargados frente a la posible descarga automatica durante el flujo?
    \item ?Las consultas van a nacer de ejemplos cerrados, de una bateria definida o de escenarios abiertos controlados?
    \item ?Que hace que el conjunto de casos sea suficientemente representativo para defender el trabajo?
\end{enumerate}

\section{Inventario del material}

\begin{table}[htbp]
\centering
\small
\begin{tabularx}{\textwidth}{L{0.22\textwidth}L{0.18\textwidth}L{0.16\textwidth}Y}
\toprule
\textbf{Elemento} & \textbf{Tipo} & \textbf{Estado} & \textbf{Papel dentro del TFM} \\
\midrule
\textit{Completar} & \textit{Datos / casos / artefactos} & \textit{Disponible / por construir} & \textit{Completar} \\
\textit{Completar} & \textit{Datos / casos / artefactos} & \textit{Disponible / por construir} & \textit{Completar} \\
\textit{Completar} & \textit{Datos / casos / artefactos} & \textit{Disponible / por construir} & \textit{Completar} \\
\bottomrule
\end{tabularx}
\caption{Inventario inicial del material de trabajo.}
\label{tab:inventario_material_trabajo}
\end{table}

\section{Casos y escenario experimental de partida}

\textit{Completar aqui el tipo de consultas que quieres estudiar: consultas directas, comparativas, tecnicas, con restricciones de formato, con necesidad de descarga, con validaciones o con reparacion. La idea es anticipar la forma de la evaluacion sin entrar todavia en las metricas formales.}

\section{Criterio de redaccion recomendado}

Cuando este capitulo quede cerrado, deberia responder con claridad a tres preguntas:

\begin{enumerate}
    \item ?Con que material concreto se construye el proyecto?
    \item ?Que tipo de casos o consultas se espera resolver?
    \item ?Por que ese punto de partida es suficiente para justificar el estudio?
\end{enumerate}
